\section{Descent along a fibration}
\label{sec:rank-one}

Throughout this section we assume Assumption~\ref{asm:log-iitaka}.

We prove Proposition~\ref{rank:fibration-case}.  Its geometric reduction
leads to a fibration whose very general fiber has logarithmic Kodaira
dimension zero.  In sufficiently divisible degrees the log pluricanonical
systems on that fiber are one-dimensional, and a relative generating form
determines an adjoint line on the base.  We prove the resulting rank-one
proposition by a secondary induction on the positive base dimension: after proving
pseudo-effectivity on the base, we either lower its dimension or lift a nef
line and its entire fixed divisor.  The last step proves generation of the
remaining nef line.

\subsection{Reduction to the rank-one case}

\begin{proposition}[The rank-one case]
\label{rank:rank-one}
Assume Assumption~\ref{asm:log-iitaka}.  Fix \(n>0\) and assume
\(\Ggood_d\) for every \(d<n\).
Let \(X\) be a smooth connected compact Kähler manifold of dimension \(n\),
let \(B\) be a rational SNC boundary, and suppose that
\(J=K_X+B\) is pseudo-effective.  Suppose there are smooth compact Kähler
manifolds \(\widetilde X,W\), a modification
\(\pi:\widetilde X\to X\), and a proper surjective holomorphic map
\(g:\widetilde X\to W\) with connected fibers such that
\[
 0<k:=\dim W<n,\qquad
 W\text{ is projective or }a(W)=0.
\]
Assume that
\(\widetilde B=\pi_*^{-1}B+\operatorname{Exc}(\pi)_{\mathrm{red}}\)
has SNC support and that, on a very general smooth fiber \(F\) of \(g\),
\[
             \kappa(F,K_F+\widetilde B_F)=0.
\]
Then \((X,B)\) satisfies \(\Ggood_n\).
\end{proposition}

\begin{proof}[Proof of Proposition~\ref{rank:fibration-case},
assuming Proposition~\ref{rank:rank-one}]
We use the reduced-exceptional boundary after every source modification.
The effective discrepancy identity in
Proposition~\ref{bir:resolution-transfer} preserves pseudo-effectivity,
and that proposition transfers the resulting decomposition back to the
original pair.  A dominant meromorphic map to
a space in class \(\mathcal C\) can be resolved on smooth compact Kähler
models; taking Stein factorization gives connected fibers.  These
operations preserve the positive dimensions of the base and the fiber.

We first suppose \(a(X)>0\).  If \(a(X)=n\), then \(X\) is Moishezon and
Kähler, hence projective, and Proposition~\ref{proj:good-models} applies.
Otherwise resolve the algebraic reduction as a fibration
\(g:X'\to W\) to a smooth projective variety of dimension \(a(X)\).
Put \(J'=K_{X'}+B'\) for the modified log adjoint.  Its restriction to
a very general smooth fiber is pseudo-effective.  Indeed, one may choose
a sequence of positive currents in
\(c_1(J')+\varepsilon_j[\omega]\), with analytic singularities and
\(\varepsilon_j\downarrow0\), and restrict all of them outside a
countable union of proper analytic subsets of the base.  By
\(\Ggood_{\dim F}\), the restricted adjoint has nonnegative Kodaira
dimension.

That dimension must be zero.  If it were positive, generic coherent
base change would give a divisible \(q\) for which
\(\mathcal E=g_*\OO_{X'}(qJ')\) has a fiber system with a
positive-dimensional image on very general fibers.  A sufficiently
large ample twist of the coherent sheaf \(\mathcal E\) is generated at
the generic point of the projective \(W\).  Evaluation would therefore
give two sections of
\(\OO_{X'}(qJ')\otimes g^*\OO_W(A)\), for one ample \(A\), whose ratio
is nonconstant on a general \(g\)-fiber.  This ratio is a meromorphic
function on \(X'\).  Every meromorphic function on \(X'\) factors
through its algebraic reduction, a contradiction.
Proposition~\ref{rank:rank-one} now applies.

Suppose next that \(a(X)=0\), and resolve the fibration in the
hypothesis as \(f:(X',B')\to Y\).  The base may be replaced by a smooth
compact Kähler model, and \(a(Y)=0\), since its meromorphic functions
pull back to \(X'\).  Write \(d=\dim X'-\dim Y>0\).
The same restriction argument and \(\Ggood_d\) show that
\[
       \kappa(F,K_F+B'_F)\geq0
\]
on very general smooth fibers.  If this integer is \(j<d\), the relative
Iitaka construction of \cite[Lemma~2.6]{CI} gives, after the same
source modifications, a factorization
\[
             X''\xrightarrow{g}W\xrightarrow{h}Y
\]
with connected fibers, \(\dim W=\dim Y+j\), and
\(\kappa(G,K_G+B''_G)=0\) on a very general \(g\)-fiber.  The cited
lemma applies to rational SNC boundaries on compact manifolds in
class \(\mathcal C\); its conclusion concerns the restricted log
adjoint, and does not require finite generation.  A Kähler model of
\(W\) has algebraic dimension zero because it is dominated by \(X''\).
Thus \(0<\dim W<n\), and Proposition~\ref{rank:rank-one} applies again.

It remains to exclude the possibility that \(K_F+B'_F\) is big.
The stable-family comparison below requires a horizontal boundary with
coefficients strictly less than one, so we first lower the boundary while
keeping the restricted adjoint big.
Let \(B'_{\mathrm{hor}}\) be the sum of the components of \(B'\)
dominating \(Y\).  There is a rational \(0<\lambda<1\) such that
\(K_F+\lambda B'_{\mathrm{hor},F}\) is big for very general \(F\).
To justify the uniform choice, exclude at once the generic base-change
and image-dimension exceptional sets for all rational \(\lambda\)
and all divisible degrees.  On a fiber outside this countable union,
openness of the big cone gives one such \(\lambda\), and a single
degree has full-dimensional image.  Generic base change makes the
same choice work very generally.  The boundary
\(\lambda B'_{\mathrm{hor}}\) is horizontal, rational SNC, and has
all coefficients strictly less than one.  The stable-family
comparison \cite[Lemma~5.2]{CI} therefore supplies a proper
generically finite cover \(Y'\to Y\), a surjection \(Y'\to R\) to a
smooth projective variety, and a projective stable family
\(\mathcal V\to R\) with positive-dimensional general fiber, such
that the main transform of \(X'\) is bimeromorphic to
\(\mathcal V\times_RY'\).

The main transform still has algebraic dimension zero.  Normalize it and
factor its proper generically finite map to \(X'\) through the normal Stein
space.  A meromorphic function descends through the birational part.
Lemma~\ref{bir:finite-norm} gives its characteristic polynomial over the
finite part, with meromorphic coefficients on \(X'\).  Those coefficients
are constant because \(a(X')=0\), so irreducibility makes the function
constant.  Normalization and modifications preserve meromorphic functions.
It follows also that \(a(Y')=0\), so the surjection from \(Y'\) to
the projective \(R\) forces \(R\) to be a point.  But then the main
transform is bimeromorphic to the product of \(Y'\) with a
positive-dimensional projective variety.  Rational functions on that
factor contradict algebraic dimension zero.  This excludes the big
case and completes the reduction.
\end{proof}

\subsection{The adjoint line on the base}

We prepare the rank-one fibration and record exactly what the relative
generating form supplies.  The construction of the Hodge line and its
positivity are those of \cite[Proposition~2.7 and Theorem~3.1]{CI}.
The relative injection and the last assertion below are consequences of
the local order calculation in that proof; they will be needed after
the base is changed.

\begin{lemma}[Prepared comparison]
\label{rank:comparison}
Let \(g_0:(X_0,B_0)\to W_0\) be a fibration with connected fibers from
a smooth compact Kähler manifold to a smooth compact manifold in
class \(\mathcal C\).  Assume \(B_0\) is a rational SNC boundary and
\(\kappa(F,K_F+B_{0,F})=0\) on very general smooth fibers.
After modifications there is a diagram
\[
\begin{tikzcd}
 X \arrow[r,"\pi"] \arrow[d,"g"'] & X_0 \arrow[d,"g_0"]\\
 W \arrow[r,"\rho"'] & W_0
\end{tikzcd}
\]
with \(X,W\) smooth compact Kähler, \(g\) a fibration with connected
fibers, and
\(B=\pi_*^{-1}B_0+\operatorname{Exc}(\pi)_{\mathrm{red}}\) SNC.
Every prime of \(X\) whose \(g\)-image has codimension at least two is
\(\pi\)-exceptional.  There are a rational SNC boundary \(T\) on \(W\),
a surjection \(p:W\to S\) with connected fibers to a smooth projective
variety, and a nef rational line \(P_S\) on \(S\).  Put
\[
        M=p^*P_S,\qquad H=K_W+T+M,\qquad J=K_X+B.
\]
If \(\dim S>0\), then \(K_S+a_0P_S\) is big for some rational
\(a_0>0\); if \(S\) is a point, \(M\sim_{\Q}0\).
There is a rational divisor \(A_*\) on \(X\) giving an actual rational
line identity
\begin{equation}
                 J\sim_{\Q}g^*H+A_*.
\label{rank:line-comparison}
\end{equation}
These data satisfy the following properties.

\begin{enumerate}
\item The horizontal part \(A_{*,\mathrm{hor}}\) is effective.
On a very general smooth fiber it is the zero divisor of a relative
generating section of a divisible multiple of \(K_F+B_F\), divided
by that multiple.

\item For each prime \(D\subset W\), let \(E\) run over the primes of
\(X\) dominating \(D\), and put \(a_E=\ord_E(g^*D)\).
Then
\begin{equation}
 (A_*)_E\geq0,\qquad
       \min_{E\to D}\frac{(A_*)_E}{a_E}=0.
\label{rank:zero-minimum}
\end{equation}
There is no assertion here about the sign at a prime whose image has
codimension at least two.

\item For all degrees \(q\) divisible by one positive integer,
multiplication by the relative generating form identifies
\[
        H^0(W,qH)\simeq H^0(X,qJ)
\]
and preserves ratios of sections.  In the same degrees, for every
proper holomorphic \(h:W\to V\), there is a natural injection
\begin{equation}
 (hg)_*\OO_X(qJ)\ \lhook\joinrel\longrightarrow\
 h_*\OO_W(qH).
\label{rank:relative-injection}
\end{equation}

\item Let \(\tau:W'\to W\) be a further smooth modification included
in another diagram with the preceding properties over the same reference
\(X_0\), using the same source-boundary convention.  Let \(H'\) be the
base line produced by that diagram.  Then, as actual rational lines,
\begin{equation}
                  H'\sim_{\Q}\tau^*H+E_H,
 \qquad E_H\geq0\ \text{\(\tau\)-exceptional}.
\label{rank:higher-base}
\end{equation}
\end{enumerate}
\end{lemma}

\begin{proof}
The flattening construction of \cite[Lemma~2.4]{CI} gives the stated
reference property, and preserves it after any finite sequence of
further base modifications.  Its source boundary is the one in
\cite[Lemma~2.1]{CI}.  We may resolve the discriminant and the
horizontal divisor of the relative generator so that all relevant
strata are smooth over the dense smooth log locus.  Proposition~2.7
of \cite{CI} then gives \(T\) and a rational parabolic Hodge line
\(M\).  A positive multiple of that line is the highest rank-one
Hodge step of a complex summand of a pure real-polarizable variation
with an integral lattice.  Theorem~3.1 of \cite{CI} gives the
factorization \(M=p^*P_S\) and the stated positivity, after a further
modification and a positive rational rescaling.  These are identities
in \(\Pic\otimes\Q\), not merely identities of classes.

Here is the order computation that gives
Equations~\eqref{rank:line-comparison}--\eqref{rank:relative-injection}.
Choose a degree \(m\) for which
\(g_*\OO_X(m(K_{X/W}+B))\) has generic rank one and all boundary
denominators are cleared.  Its reflexive hull is a line on the smooth
\(W\).  Let \(s\) be a meromorphic relative generating form in a
local frame of this line.  At the generic point of a prime \(D\),
choose a nonvanishing ordinary base volume \(\omega_W\), and set
\[
 l_E=\frac1m\ord_E(s\wedge g^*\omega_W^m),\qquad
 \delta_E=B_E,\qquad
 \alpha_D=\min_{E\to D}\frac{l_E+\delta_E}{a_E}.
\]
The relative-times-base identification in this formula is taken in
the \(m\)-th canonical tensor power.  The source is resolved over
the generic point of \(D\), so all the primes needed for the
divisorial test occur in this minimum.  In the notation of the
proof of \cite[Proposition~2.7]{CI}, the parabolic order is
\[
 \beta_D=\min_{E\to D}\frac{l_E+1-a_E}{a_E},
                    \qquad T_D=\alpha_D-\beta_D.
\]
Thus the corresponding local frame of \(T+M\) has order
\(\alpha_D\).  Comparing its pullback with the log adjoint form
defines \(A_*\), with
\begin{equation}
                (A_*)_E=l_E+\delta_E-a_E\alpha_D .
\label{rank:vertical-order}
\end{equation}
Changing the generating form multiplies it by a meromorphic base
function.  This changes both \(\alpha_D\) and \(\beta_D\) by its
normalized base order and leaves the comparison invariant.  The local
comparisons therefore give the rational divisor and the line identity
globally.  Equation~\eqref{rank:vertical-order} proves
Equation~\eqref{rank:zero-minimum}.  Horizontally the comparison is
precisely the zero divisor of the fiberwise log section, so it is
effective and has the asserted restriction.

The same order computation proves the relative injection.
In a divisible degree \(q\), a meromorphic base coefficient
\(\varphi\) gives a regular base section at \(D\) exactly when
\[
                     \ord_D(\varphi)+q\alpha_D\geq0.
\]
Its corresponding upstairs orders at the primes \(E\to D\) are
\[
             a_E\ord_D(\varphi)+q(l_E+\delta_E).
\]
Their simultaneous nonnegativity is equivalent to the preceding
inequality.  Conversely, an upstairs section restricts to a multiple
of the generator on general connected fibers, so its ratio with the
generator descends meromorphically to \(W\).  The order inequalities
make it regular outside codimension two on \(W\), and normality
extends it.  Consequently
\[
                 g_*\OO_X(qJ)\ \lhook\joinrel\longrightarrow\
                 \OO_W(qH).
\]
Applying the left exact functor \(h_*\) proves
Equation~\eqref{rank:relative-injection}.  The reverse global
comparison follows from the reference property: its only initially
untested poles are on \(g\)-contracted primes, all exceptional over
\(X_0\); pushing to \(X_0\), extending in codimension two, and pulling
back with the reduced-exceptional boundary removes them.  This is
the global comparison in \cite[Proposition~2.7]{CI}.

Finally consider \(\tau:W'\to W\).  The parabolic rational line pulls
back under SNC modifications, as in the construction preceding
\cite[Theorem~3.1]{CI}.  The coefficient \(T_D\) at an old strict
transform is unchanged.  Indeed the sheaf of regular relative
adjoint forms inside the common meromorphic generator line is
unchanged over the generic point of \(D\) by the log modification
formula.  Equivalently, the old inequalities above already test all
vertical primes there; in SNC coordinates logarithmic pullback
preserves those inequalities, so new source exceptional primes
impose no stronger one.  This keeps \(\alpha_D\) unchanged in
compatible frames, while parabolic pullback keeps \(\beta_D\)
unchanged.  Over a new \(\tau\)-exceptional prime every source prime
is exceptional over \(X_0\): a pre-existing one had image of
codimension at least two on \(W\), and a new one is exceptional
by construction.  All therefore have boundary coefficient one.
The formula for \(T_D=\alpha_D-\beta_D\) then gives coefficient one
on that new base prime.  Thus
\[
      T'=\tau_*^{-1}T+\operatorname{Exc}(\tau)_{\mathrm{red}}.
\]
The usual log discrepancy formula for the SNC pair \((W,T)\)
now gives \(K_{W'}+T'=\tau^*(K_W+T)+E_H\), with \(E_H\)
effective and exceptional.  Adding the pulled-back parabolic line
proves Equation~\eqref{rank:higher-base}.
\end{proof}

\begin{lemma}[Pseudo-effectivity of the base line]
\label{rank:base-psef}
In the setting of Lemma~\ref{rank:comparison}, suppose
\(\dim X=n\), \(0<\dim W<n\), \(J\) is pseudo-effective, and
\(\Ggood_d\) holds for \(d<n\).  Then
\[
        A_{*,\mathrm{hor}}\leq N(J)
        \quad\text{as divisors},\qquad H\ \text{is pseudo-effective}.
\]
\end{lemma}

\begin{proof}
On a very general smooth fiber \(F\), use the model supplied by
\(\Ggood_{\dim F}\) to write
\(\mu_F^*J_F\sim_{\Q}P_F+N(\mu_F^*J_F)\).
Lemma~\ref{bir:sections} identifies the divisible section spaces with
those of \(P_F\), so \(\kappa(P_F)=\kappa(F,J_F)=0\); a semiample line of Kodaira
dimension zero is torsion.  Proposition~\ref{bir:finite-descent},
applied to the modification \(\mu_F\) with zero added divisor, gives
\(J_F\sim_{\Q}N(J_F)\) on \(F\), with \(N(J_F)\) rational.
Lemma~\ref{bir:sections} then identifies the normalized zero divisor of
the prepared relative generator with \(N(J_F)\).  Thus
\[
        A_{*,\mathrm{hor}}|_F=N(J_F).
\]

We compare these fiber multiplicities with those on \(X\).
For each of the finitely many components \(E\) of
\(A_{*,\mathrm{hor}}\), choose a countable sequence of
small-Kähler-perturbation currents with analytic singularities whose
generic orders at \(E\) approach \(\nu_E(J)\).  Choose a common very
general \(F\) for these sequences and these components.  The currents restrict
positively after the same perturbations, and analytic singularities
ensure that their orders along the components of \(E|_F\) are
their generic orders along \(E\).  Each restricted order is at
least the corresponding perturbed minimal multiplicity on \(F\).
Passing to the limit gives
\((A_*)_E\leq\nu_E(J)\).  This proves
\begin{equation}
                       A_{*,\mathrm{hor}}\leq N(J).
\label{rank:horizontal-negative}
\end{equation}

Choose a positive curvature current for the rational line \(J\).
Every such current contains its divisorial negative part, hence
subtracting \([A_{*,\mathrm{hor}}]\) leaves a positive current.
On the inverse image of a dense smooth open \(W^\circ\), with all
vertical data removed, Equation~\eqref{rank:line-comparison}
identifies the resulting singular metric with one on \(g^*H\).
In a frame pulled back from \(W^\circ\) its weight is plurisubharmonic.
It is constant on each compact connected smooth fiber, by the
maximum principle (with the value \(-\infty\) allowed).  Local
holomorphic sections of the submersion show that these values form
a plurisubharmonic weight on \(H|_{W^\circ}\).

This weight extends across the missing divisors.  At the generic
point of any such divisor \(D\), choose \(E\to D\) for which
\((A_*)_E=0\), using Equation~\eqref{rank:zero-minimum}.
At a general point of \(E\) away from the other comparison divisors,
the comparison after subtracting the horizontal part has neither
a zero nor a pole.  We can choose local coordinates, with the remaining
vertical coordinates denoted by \(w\), in which
\[
 g(u,z_2,\ldots,z_k,w)=(u^{a_E},z_2,\ldots,z_k).
\]
Indeed \(g|_E\) has maximal tangential rank generically and a local
equation of \(g^*D\) is a unit times \(u^{a_E}\); the unit has a
local \(a_E\)-th root.  A smaller such chart covers a neighborhood
of the base point.  The upstairs plurisubharmonic weight is locally
bounded above there, and the comparison unit is bounded.  Hence
the descended weight is locally bounded above near the generic
point of \(D\).  The removable singularity theorem for
plurisubharmonic functions extends it across \(D\) away from
codimension two, and the Hartogs extension for plurisubharmonic
functions extends it across the remaining analytic subset.
The extensions respect the line transitions because they agree on
the dense open.  They give a positive singular metric on \(H\),
which proves pseudo-effectivity.
\end{proof}

\begin{lemma}[A base of algebraic dimension zero]
\label{rank:zero-base}
Under the assumptions of Lemma~\ref{rank:base-psef}, if \(a(W)=0\),
then, on \(W\),
\begin{equation}
                        H\sim_{\Q}E_W=N(H)
                        \quad\text{with }E_W\geq0.
\label{rank:zero-base-decomposition}
\end{equation}
\end{lemma}

\begin{proof}
The projective quotient \(S\) in Lemma~\ref{rank:comparison} is a
point: otherwise its rational functions pull back nontrivially to
\(W\).  Thus \(M\sim_{\Q}0\) and \(H\sim_{\Q}K_W+T\).
By Lemma~\ref{rank:base-psef} and \(\Ggood_{\dim W}\), its pullback
to a smooth higher model is a semiample rational line plus its
negative divisor.  Every semiample line on a space of algebraic
dimension zero is torsion, since its generated system has
zero-dimensional image.  Proposition~\ref{bir:finite-descent}, applied
to this modification with zero added divisor, therefore gives
Equation~\eqref{rank:zero-base-decomposition}, including rationality
of the divisor.
\end{proof}

\subsection{Projective programs under the assumption}
\label{rank:projective-program-input}

The argument for a projective base uses a terminating program to reach a nef
adjoint and, after decreasing the nef data, a program ending in a Mori fiber
space. The conditional input is termination in the pseudo-effective case.
We first state the exact program result and then derive that case from the
projective good models in Proposition~\ref{proj:good-models}.

\begin{proposition}[A conditional generalized program]
\label{proj:generalized-mmp}
Assume Assumption~\ref{asm:log-iitaka}. Let
\((V,B+\boldsymbol M)\) be a projective generalized dlt rational pair,
where \(V\) is \(\Q\)-factorial, \(B\geq0\), and the nef b-divisor
\(\boldsymbol M\) is determined by a nef rational Cartier divisor on a
projective birational model. Put \(D=K_V+B+M_V\).

There is a choice of \(D\)-MMP with ample scaling which terminates. If
\(D\) is pseudo-effective, its endpoint \(V_m\) has nef adjoint \(D_m\).
If \(D\) is not pseudo-effective, its endpoint has a Mori fiber contraction
to a projective variety of smaller dimension. The birational steps and their
endpoints stay \(\Q\)-factorial generalized dlt, and the forward birational
map extracts no divisors.
In the pseudo-effective case, on a common resolution one has an actual
rational divisor identity
\[
 u^*D=v^*D_m+F,\qquad F\geq0\text{ and }v\text{-exceptional}.
\]
\end{proposition}

For the pseudo-effective alternative, we first derive the relative weak
Zariski decompositions used by the minimal-model existence theorem.

For a projective morphism, first replace its image by its normal Stein
factor. We use the relative convention of \cite[Section~2]{TsakanikasXie}:
a divisor is pseudo-effective over the base when its restriction to a very
general fiber of this surjective morphism is pseudo-effective. An NQC weak Zariski decomposition of \(D\) over the
base is a birational equality \(h^*D\equiv_Z P+N\), where \(P\) is a
nonnegative real combination of relatively nef rational Cartier divisors
and \(N\geq0\). The decomposition below has rational Cartier parts.

\begin{lemma}[Relative canonical decomposition]\label{proj:relative-wzd}
Assume Assumption~\ref{asm:log-iitaka}. Let \(V\) be a smooth
quasi-projective complex variety and let \(f:V\to Z\) be projective, with
\(Z\) normal and quasi-projective. If \(K_V\) is pseudo-effective over
\(Z\), then on a normal variety \(Y\) with a projective birational morphism \(h:Y\to V\) there
are rational Cartier divisors \(P,N\) such that
\[
 h^*K_V\equiv_Z P+N,\qquad P\text{ nef over }Z,\qquad N\geq0.
\]
\end{lemma}

\begin{proof}
We use the construction in the proof of
\cite[Proposition~3.3]{External037}. Its absolute input is a nef-plus-
effective rational decomposition for the canonical divisor of a smooth
projective variety with pseudo-effective canonical class. This is supplied
here by Proposition~\ref{proj:good-models} and its common-resolution
comparison.

Replace the image of \(f\) by its normal Stein factor. This does not change
relative numerical classes or contracted curves, and now \(f\) is
surjective with connected fibers. Its very general fiber is smooth
projective with pseudo-effective canonical divisor, hence is non-uniruled
by \cite[Theorem~0.2 and Corollary~0.3]{BDPP}. Compactify the projective
morphism and resolve away from \(V\). We obtain
\(\bar f:\bar V\to\bar Z\), with \(\bar V\) smooth projective and
\(\bar Z\) normal projective, unchanged over \(Z\). If \(\dim Z=0\),
\(V\) is already projective and the absolute input proves the lemma.

Suppose \(\dim Z>0\), and let \(r:\widetilde Z\to\bar Z\) be a smooth
projective resolution. Choose a sufficiently positive very ample divisor
\(A_Z\) on \(\bar Z\) and a general \(B\in|2A_Z|\). The pulled-back systems
on both \(\widetilde Z\) and \(\bar V\) are base point free. Bertini makes
\(r^*B\) and \(\bar f^*B\) smooth nonempty reduced divisors. Their double
covers
\[
 \pi_Z:Z^\sharp\to\widetilde Z,
 \qquad \pi:V^\sharp\to\bar V
\]
are smooth integral projective varieties. The canonical formulas are
\[
 K_{Z^\sharp}\sim_{\Q}\pi_Z^*(K_{\widetilde Z}+r^*A_Z),
 \qquad
 K_{V^\sharp}\sim_{\Q}\pi^*(K_{\bar V}+\bar f^*A_Z).
\]
The first line is big when \(A_Z\) is sufficiently positive. The rational
map \(V^\sharp\dashrightarrow Z^\sharp\) has the same very general fibers
as \(f\). A covering family of rational curves on \(V^\sharp\) would
either dominate a covering family of rational curves on \(Z^\sharp\), or
cover its very general fibers by vertical rational curves. Both are
impossible. Thus \(V^\sharp\) is non-uniruled and \(K_{V^\sharp}\) is
pseudo-effective by BDPP.

Apply the absolute decomposition to \(V^\sharp\). Moving the rational
principal difference in the canonical formula into the nef part gives an
actual rational divisor decomposition for a pullback of
\(\pi^*(K_{\bar V}+\bar f^*A_Z)\). Take a common equivariant resolution of
that model and its conjugate under the covering involution. On this smooth
projective resolution \(g:T\to V^\sharp\), average the two decompositions.
We obtain
\[
 g^*\pi^*(K_{\bar V}+\bar f^*A_Z)=\bar P+\bar N,
 \qquad \bar P\text{ nef},\quad \bar N\geq0,
\]
where \(\bar P\) and \(\bar N\) are invariant rational Cartier divisors.

Let \(q:T\to\bar Y\) be the finite quotient by the involution. The invariant
composite \(\pi g\) induces a projective birational morphism
\(\bar h:\bar Y\to\bar V\), with \(\bar Y\) normal and projective. An invariant
rational divisor descends by dividing its coefficient at an upstairs prime
by the ramification index. The descended divisor is rational Cartier when
the original divisor is: after clearing denominators, choose one Cartier
equation on a semilocal neighborhood of the finite orbit above a point
(a Cartier divisor is principal on a semilocal ring), and multiply all its
translates. The product is invariant and has divisor
equal to the group order times the original divisor. It descends to a local
equation for a multiple of the downstairs divisor. This argument includes
stabilizers and ramification primes.

Accordingly \(\bar P=q^*P\) and \(\bar N=q^*N\) for rational Cartier
divisors on \(\bar Y\). Lifting curves through the finite map shows that
\(P\) is nef, and coefficientwise descent shows that \(N\) is effective.
Finite pullback is injective on rational divisors, so the displayed identity
descends. Restrict it to \(Y=\bar h^{-1}(V)\). The term pulled
back from \(A_Z\) is numerically trivial over \(Z\), giving
\(h^*K_V\equiv_ZP+N\), as required.
\end{proof}

\begin{proof}[Proof of Proposition~\ref{proj:generalized-mmp}]
Lemma~\ref{proj:relative-wzd} and
\cite[Theorem~5.2]{TsakanikasXie}, applied to the generalized pair
\((U/Z,0+0)\) with \(U\) smooth, prove relative minimal-model existence for smooth varieties
by dimension induction. The hypothesis of that theorem is relative smooth
minimal-model existence one dimension lower; the lemma supplies its other
hypothesis, an NQC weak Zariski decomposition. The induction begins in
dimension zero. Then \cite[Theorem~5.4]{TsakanikasXie} supplies minimal
models for every pseudo-effective NQC generalized lc pair in the relevant
dimension. Their Theorem~2.7 gives existence in the Birkar--Shokurov sense.

The nef data in the statement are NQC, because a nef rational Cartier
divisor is a positive multiple of a nef Cartier divisor. In the
pseudo-effective case the preceding paragraph therefore gives the
Birkar--Shokurov model required by
\cite[Theorem~4.2]{TsakanikasXie}. In the other case, the non-pseudo-effective
alternative of that theorem applies directly. To meet its scaling
hypothesis, take an effective sufficiently positive ample rational divisor
\(A\) whose general components have sufficiently small coefficients. On
a fixed log resolution carrying \(\boldsymbol M\), these components are
transverse to the boundary and preserve generalized log canonicity, while
their ample class can be chosen large enough that \(D+A\) is nef. The
theorem supplies a terminating MMP starting on \(V\).

We verify the step types inductively.  Suppose \(V_i\) is
\(\Q\)-factorial generalized dlt.  By
\cite[Remark~2.15]{TsakanikasXie}, each birational step has the form
\[
 V_i\xrightarrow{g_i}Z_i\xleftarrow{h_i}V_{i+1},
\]
where \(g_i\) contracts a negative extremal ray and \(h_i\) is small and projective.
Suppose \(g_i\) is divisorial, and write \(R\) for its ray.  An exceptional
prime \(E\) has \(E\cdot R<0\): otherwise it would be effective,
exceptional, and \(g_i\)-nef, contrary to negativity.  There is only one
exceptional prime.  Indeed, a combination of two distinct exceptional
primes can be chosen to have degree zero on \(R\); negativity applied to
both signs of that combination would make it zero.

We show directly that \(Z_i\) is \(\Q\)-factorial.  For a prime divisor
\(T\) on \(Z_i\), let \(\widetilde T\) be its strict transform, choose a
curve \(C\) spanning \(R\), and put
\[
                 a=-\frac{\widetilde T\cdot C}{E\cdot C}\in\Q.
\]
Choose a positive integer \(m\) for which \(m(\widetilde T+aE)\) is
Cartier.  Its line has degree zero on \(R\), so
\cite[Theorem~1.5]{Xie2022} gives an actual Cartier line \(\mathcal L_T\)
on \(Z_i\) and an isomorphism
\[
 \OO_{V_i}\bigl(m(\widetilde T+aE)\bigr)\simeq g_i^*\mathcal L_T.
\]
On an open set trivializing \(\mathcal L_T\), the Cartier divisor
\(m(\widetilde T+aE)\) is principal upstairs.  Pushing this principal
divisor identity down through the birational \(g_i\) gives a principal
divisor \(mT\) on that open set.
Thus \(T\) is rational Cartier.  It follows that \(h_i\) is an isomorphism:
the pushforward of an \(h_i\)-ample Cartier divisor is rational Cartier
on \(Z_i\), and its pullback is the original divisor because \(h_i\)
is small.  It has degree zero on every contracted curve, so relative
ampleness rules out positive-dimensional fibers; a finite birational
map to the normal \(Z_i\) is an isomorphism.  If \(g_i\) is small,
the relatively ample canonical model \(h_i\) is the flip of
\cite[Definition~3.2]{HanLi}.  Thus the chosen birational steps are
divisorial contractions and flips.  The underlying variety of a
generalized dlt pair is klt, and these steps preserve the
\(\Q\)-factorial generalized dlt category
\cite[Remark~2.3 and Lemma~3.7]{HanLi}.  They introduce no prime
divisors on the new models. For completeness, on a common resolution
of one step \(V_i\dashrightarrow V_{i+1}\) carrying \(\boldsymbol M\),
taking the difference of the two generalized discrepancy formulas cancels
their common nef divisor. The resulting difference of the adjoint pullbacks is exceptional over the new
model and anti-nef over it. The negativity lemma makes this difference
effective. Composing these actual rational divisor comparisons gives
\(u^*D=v^*D_m+F\) with \(F\) effective and exceptional over the endpoint.
The endpoints in the two cases are the nef and Mori fiber endpoints supplied
by the terminating program theorem.
\end{proof}


\subsection{Reducing a projective base}

The nef line \(M\) supplied by the Hodge construction need not be
semiample.  A first projective program replaces \(H\) by a nef transform.
Positive Iitaka dimension gives either a smaller base or a big line.  In
nonpositive Iitaka dimension, rationally trivial nef data give a trivial
line.  In the remaining case,
decreasing \(M\) leads to a second program that lowers the base dimension
while preserving the first nef line.

\begin{lemma}[Projective base reduction]
\label{rank:base-reduction}
Assume Assumption~\ref{asm:log-iitaka} and the setting of
Lemma~\ref{rank:base-psef}, with \(W\) projective of dimension \(k\).
Then one of the following holds.
\begin{enumerate}
\item After a smooth source modification with the reduced-exceptional
boundary, there is a fibration to a smooth projective variety \(V\)
with \(0<\dim V<k\) and logarithmic Kodaira dimension zero on very
general smooth fibers.
\item There are a normal \(\Q\)-factorial projective variety \(W_m\),
a birational contraction \(W\dashrightarrow W_m\) extracting no
divisors, and a nef rational line \(H_m\) on \(W_m\), either big or
rationally trivial, such that on a common smooth resolution
\(\mu:W'\to W\), \(\nu:W'\to W_m\),
\begin{equation}
             \mu^*H\sim_{\Q}\nu^*H_m+E,
             \qquad E\geq0\ \text{\(\nu\)-exceptional}.
\label{rank:base-endpoint}
\end{equation}
\end{enumerate}
Here all the line comparisons are actual rational line identities.
\end{lemma}

\begin{proof}
The smooth pair with boundary \(T\) and nef data \(\boldsymbol M\)
determined by \(M\) on \(W\) is a rational \(\Q\)-factorial
generalized dlt pair.  Its adjoint \(H\) is pseudo-effective by
Lemma~\ref{rank:base-psef}.  Proposition~\ref{proj:generalized-mmp}
gives a chosen terminating program with scaling to a
\(\Q\)-factorial generalized dlt model \(W_{\mathrm{nef}}\) with nef transformed
adjoint
\[
 H_{\mathrm{nef}}
     =K_{W_{\mathrm{nef}}}+T_{\mathrm{nef}}+M_{\mathrm{nef}}.
\]
Here \(M_{\mathrm{nef}}\) is the trace of the nef b-divisor whose nef
representative is \(M\) on \(W\).  On a common smooth resolution
\(\mu:W'\to W\), \(\sigma:W'\to W_{\mathrm{nef}}\), the actual comparison is
\begin{equation}
 \mu^*H\sim_{\Q}\sigma^*H_{\mathrm{nef}}+E_{\mathrm{nef}},
 \qquad E_{\mathrm{nef}}\geq0\ \text{\(\sigma\)-exceptional}.
\label{rank:first-nef-endpoint}
\end{equation}
The program extracts no divisors.  We use this particular model also when
the nef data are rationally trivial.

If \(\kappa(H)>0\), Equation~\eqref{rank:first-nef-endpoint} preserves
the divisible section spaces.  For \(\kappa(H)=k\) the nef
\(H_{\mathrm{nef}}\) is big, and the second alternative holds with
\(W_m=W_{\mathrm{nef}}\) and \(H_m=H_{\mathrm{nef}}\).  If
\(0<\kappa(H)<k\), the exact global comparison in
Lemma~\ref{rank:comparison} gives \(\kappa(J)=\kappa(H)\).
Take the Iitaka fibration of \(J\) on the source.  Its projective
image, followed by Stein factorization and a projective resolution,
has dimension \(\kappa(H)\).  The log adjoint on very general
fibers has Kodaira dimension zero by the relative Iitaka construction
\cite[Lemma~2.6]{CI}, including its persistence under the
reduced-exceptional source convention.  This is the first alternative.

Suppose \(\kappa(H)\leq0\) and \(M\sim_{\Q}0\).  We may choose zero
nef data in this rational equivalence class.  The same generalized
dlt model \(W_{\mathrm{nef}}\) is then an ordinary dlt model and
\(H_{\mathrm{nef}}\sim_{\Q}K_{W_{\mathrm{nef}}}+T_{\mathrm{nef}}\).
Its adjoint is nef, so
Proposition~\ref{proj:good-models} makes \(H_{\mathrm{nef}}\) semiample as an
actual rational line.  Since
\(\kappa(H_{\mathrm{nef}})=\kappa(H)\leq0\), a generated multiple has
constant image and is trivial.  Thus the second alternative holds with
\(W_m=W_{\mathrm{nef}}\) and \(H_m=H_{\mathrm{nef}}\sim_{\Q}0\).

It remains to suppose
\begin{equation}
                    \kappa(H)\leq0,\qquad M\not\sim_{\Q}0.
\label{rank:nontrivial-nef-data}
\end{equation}
Here \(\dim S>0\).  For a rational \(c\in(0,1)\), put
\(H_c=K_W+T+cM\).  We claim
\begin{equation}
                   H_c\ \text{is not pseudo-effective}
                   \quad(0<c<1,\ c\in\Q).
\label{rank:lowered-nonpsef}
\end{equation}
Suppose otherwise for one \(c\), and choose an ample rational line \(U\)
on \(S\).  We will find \(\delta>0\) for which \(H-\delta p^*U\) has an
effective representative, forcing \(\kappa(H)\geq\dim S>0\).
Choose a rational \(a>\max\{1,a_0\}\) and a small
rational \(\lambda>0\) such that \(K_S+aP_S-\lambda U\) is big.
This is possible because \(P_S\) is nef and
\(K_S+a_0P_S\) is big.  Choose
\[
                   0<\eta<\lambda\frac{1-c}{a-1}
                   \quad\text{rational}.
\]
The rational line \(cP_S+\eta U\) is ample.  A sufficiently high
general member divided by its degree pulls back to a rational
boundary transverse to \(T\) with subunit coefficient.  Applying
the nonvanishing part of Proposition~\ref{proj:good-models} to
this ordinary projective log pair gives
\[
                 E_c\sim_{\Q}H_c+\eta p^*U,\qquad E_c\geq0.
\]
A nonzero section defining \(E_c\) restricts nontrivially to very
general \(p\)-fibers in a fixed divisible degree.  The hypotheses
of the weak-effectivity statement \cite[Lemma~3.3]{CI} are now
satisfied: the source is smooth in class \(\mathcal C\), \(T\)
is rational SNC, the base \(S\) is smooth projective, and a
relative log system is nonzero.  Applied to the big line
\(K_S+aP_S-\lambda U\), it gives
\[
                 E_a\sim_{\Q}K_W+T+aM-\lambda p^*U,\qquad E_a\geq0.
\]
For
\[
 \theta=\frac{1-c}{a-c},\qquad
 \delta=\theta\lambda-(1-\theta)\eta>0,
\]
the rational interpolation is
\[
          (1-\theta)E_c+\theta E_a
                    \sim_{\Q}H-\delta p^*U.
\]
Multiplying pullbacks of sections of \(U\) by a section of the
effective left-hand side yields
\(\kappa(H)\geq\kappa(p^*U)=\dim S>0\).
This contradicts Equation~\eqref{rank:nontrivial-nef-data} and
proves Equation~\eqref{rank:lowered-nonpsef}.

We next run a program to fiber type while descending the actual nef line
\(H_{\mathrm{nef}}\) through every contraction.  Choose the common
resolution in Equation~\eqref{rank:first-nef-endpoint} to determine the nef
data.  The negativity lemma gives
\[
                F_M:=\sigma^*M_{\mathrm{nef}}-\mu^*M\geq0
                \quad\text{\(\sigma\)-exceptional}.
\]
Indeed its negative is \(\sigma\)-nef and exceptional.  For every
rational \(0<e<1\), Equation~\eqref{rank:first-nef-endpoint} gives
\begin{equation}
 \mu^*H_{1-e}\sim_{\Q}
       \sigma^*(H_{\mathrm{nef}}-eM_{\mathrm{nef}})
       +E_{\mathrm{nef}}+eF_M .
\label{rank:lowered-comparison}
\end{equation}
If \(H_{\mathrm{nef}}-eM_{\mathrm{nef}}\) were pseudo-effective, this
equality and effectivity of the last two terms would make
\(\mu^*H_{1-e}\), and hence \(H_{1-e}\), pseudo-effective, contrary to
Equation~\eqref{rank:lowered-nonpsef}.

Fix rational \(0<\varepsilon<1\), choose an integer \(r>0\) with
\(rH_{\mathrm{nef}}\) Cartier, and choose an integer \(d_0>2kr\).
This coefficient will force every ray of the second program to have zero
degree for the descended nef line.  Decreasing the nef data to
\((1-\varepsilon)\boldsymbol M\) preserves generalized dlt
singularities: on the displayed resolution the crepant boundary
changes by \(-\varepsilon F_M\), so discrepancies do not decrease.
Add \(d_0H_{\mathrm{nef}}\), determined as nef data on
\(W_{\mathrm{nef}}\).  It changes no discrepancies, and the resulting
adjoint is
\[
 (1+d_0)H_{\mathrm{nef}}-\varepsilon M_{\mathrm{nef}}
 =(1+d_0)\left(
       H_{\mathrm{nef}}-\frac{\varepsilon}{1+d_0}M_{\mathrm{nef}}\right).
\]
It is not pseudo-effective by the preceding paragraph.  The
non-pseudo-effective case of Proposition~\ref{proj:generalized-mmp}
therefore gives a chosen terminating program to a Mori fiber
contraction in the \(\Q\)-factorial generalized dlt category.

At any stage of this second program, let \(H_i\) be the descended rational
line, assuming inductively that it is nef and \(rH_i\) is Cartier, and write
\[
              G_i=K_{W_i}+T_i+(1-\varepsilon)M_i.
\]
An extremal ray negative for \(G_i+d_0H_i\) is negative for \(G_i\).
The pair with adjoint \(G_i\) is still generalized dlt: the added
nef b-divisor \(d_0H_i\) descends to the current model, so removing
it changes no discrepancy.  The generalized length bound
\cite[Proposition~3.13]{HanLi} gives a curve \(C\)
on that ray with \(0<-G_i\cdot C\leq2k\).  If \(H_i\cdot C>0\),
Cartier integrality would give
\[
               \frac{d_0}{r}\leq d_0H_i\cdot C
                       <-G_i\cdot C\leq2k,
\]
a contradiction.  Thus \(H_i\cdot C=0\).
The exact contraction statement \cite[Theorem~1.5]{Xie2022}
descends the Cartier line \(rH_i\) along this contraction as an
actual Cartier line.  In a flip we pull that same line to the
flipped side.  The descended line is nef: every curve on the
contraction base has a curve above it mapping with positive degree,
so its degree is nonnegative; pullback preserves nefness.
Consequently \(r\) and nefness persist, and the argument applies
inductively to all steps and to the final contraction.  Each step
is crepant for \(H_i\).

Let \(W_m\) be the last birational model of this second program and let
\(H_m\) be the descended line on it.  Crepancy for the \(H_i\), together
with the absence of extraction, turns
Equation~\eqref{rank:first-nef-endpoint} into
Equation~\eqref{rank:base-endpoint} on a common resolution, with an
effective divisor exceptional over \(W_m\).
The final contraction \(h:W_m\to V\) has connected fibers and,
by the same exact descent,
\[
                       H_m\sim_{\Q}h^*H_V
\]
for a rational line \(H_V\) on the normal projective \(V\).
If \(V\) is a point, the second alternative holds with
\(H_m\sim_{\Q}0\).

Assume \(\dim V>0\).  To obtain the first alternative, it remains to show
that the log adjoint on the new very general fibers has Kodaira dimension
zero.  Resolve the base program by
\(\mu:W''\to W\) and \(\nu:W''\to W_m\), and prepare the source
over \(W''\), calling it \(g'':(X'',B'')\to W''\).  The diagram
we use is
\[
\begin{tikzcd}
 X'' \arrow[r,"g''"] & W'' \arrow[r,"\nu"] \arrow[d,"\mu"'] &
 W_m \arrow[r,"h"] & V\\
 & W
\end{tikzcd}
\]
with composite \(f= h\nu g''\).
Equation~\eqref{rank:higher-base} and
Equation~\eqref{rank:base-endpoint} give
\[
                  H''\sim_{\Q}\nu^*H_m+E'',
                  \qquad E''\geq0\ \text{\(\nu\)-exceptional}.
\]
To see the exceptionality, the error in
Equation~\eqref{rank:base-endpoint} is already exceptional over
\(W_m\), and every \(\mu\)-exceptional prime is exceptional
over \(W_m\) as well: otherwise \(W_m\) would extract a divisor
over \(W\).
For every degree divisible by one fixed integer,
Equation~\eqref{rank:relative-injection}, exceptional descent,
and projection formula now yield
\begin{align*}
 f_*\OO_{X''}(q(K_{X''}+B''))
 &\lhook\joinrel\longrightarrow
       (h\nu)_*\OO_{W''}(qH'')\\
 &=h_*\OO_{W_m}(qH_m)=\OO_V(qH_V).
\end{align*}
Here \(\nu_*\OO_{W''}(qE'')=\OO_{W_m}\) by normality and
\(h_*\OO_{W_m}=\OO_V\).  Generic coherent base change,
simultaneously for the countably many divisible degrees, bounds
the dimension of every such system on a very general \(f\)-fiber
by one.  The log adjoint on that fiber is pseudo-effective by
restriction of a countable sequence of perturbed currents with
analytic singularities.  Its dimension is less than \(n\), so
the lower-dimensional \(\Ggood\) supplies a nonzero section in
some divisible degree.  Taking a common multiple shows that the
fiber log Kodaira dimension is zero.  All maps in the composite
have connected fibers; resolving \(V\) and the source keeps this
property and gives a smooth projective base of dimension
\(0<\dim V<k\).  This is the first alternative.
\end{proof}

\subsection{The negative divisor upstairs}

The next lemma lifts either base decomposition and determines the whole
negative divisor on the source, including primes over subsets of
codimension at least two in the base that the local comparison did not test.

\begin{lemma}[Lifting the decomposition]
\label{rank:lift}
Assume the setting of Lemma~\ref{rank:base-psef}.  Suppose one of
the following additional conditions holds:
\begin{enumerate}
\item \(a(W)=0\) and \(H\sim_{\Q}E_W=N(H)\), with \(E_W\geq0\);
\item there is a birational morphism \(\nu:W\to W_m\) to a normal
\(\Q\)-factorial projective variety and an identity
\[
              H\sim_{\Q}\nu^*H_m+E_W,
              \qquad E_W\geq0\ \text{\(\nu\)-exceptional},
\]
where \(H_m\) is a nef rational line, either big or rationally
trivial.
\end{enumerate}
Then there is an effective rational divisor \(A\) on \(X\) and an actual
rational line identity
\begin{equation}
                J\sim_{\Q}P_0+A,\qquad A=N(J)\geq0,
\label{rank:upstairs-decomposition}
\end{equation}
where \(P_0\sim_{\Q}0\) in the first case and in the rationally
trivial part of the second case, and
\(P_0=(\nu g)^*H_m\) in the big part of the second case.
\end{lemma}

\begin{proof}
Use the chosen base identity in Equation~\eqref{rank:line-comparison}.
In a rationally trivial case fix a rational trivialization of the positive
line and put \(P_0=0\); in the big case put \(P_0=(\nu g)^*H_m\).  Define
the rational divisor
\[
                         A:=A_*+g^*E_W.
\]
The composed actual line identities give \(J\sim_{\Q}P_0+A\).
The divisor \(A\) may at first have signed coefficients.  We first prove
\(A\geq0\).

In a rationally trivial case a divisible multiple of the base line
has a section whose normalized zero divisor is exactly \(E_W\).
Its section under the exact global comparison has normalized zero
divisor \(A\).  Regularity of this section proves \(A\geq0\),
also on the initially untested primes.
In the big case choose a Kodaira decomposition
\[
                     H_m\sim_{\Q}A_0+G_0,
                     \qquad A_0\text{ ample},\quad G_0\geq0.
\]
For rational \(0<t<1\),
\[
 H_m\sim_{\Q}
       \underbrace{(1-t)H_m+tA_0}_{\text{ample}}+tG_0.
\]
Fix a prime \(\Gamma\subset X\).  A sufficiently divisible
section of the ample summand can be chosen not to vanish at the
generic point of \((\nu g)(\Gamma)\); its pullback has zero order
at \(\Gamma\).  Multiplication by the section of \(tG_0\) and
then the global comparison gives a regular upstairs section.
Its normalized order is
\[
               A_\Gamma+
                    t\,\ord_\Gamma((\nu g)^*G_0)\geq0.
\]
Letting \(t\downarrow0\) proves \(A_\Gamma\geq0\).  This applies
to every prime in \(A\).

Nefness of \(P_0\) now implies \(N(J)\leq A\).
Set \(Q=A-N(J)\geq0\), an effective real divisor.  By
Equation~\eqref{rank:horizontal-negative}, \(Q\) is vertical.
Moreover
\[
                    \{P_0+Q\}=\{J\}-\{N(J)\}
\]
is modified nef.  We prove \(Q=0\).

Let \(b:X\to Y\) denote \(g:X\to W\) in the first case and
\(\nu g:X\to W_m\) in the second.  Thus \(\dim Y=k\), and \(b\)
has connected fibers.  If \(Q\ne0\), let
\[
       l=\max\{\dim b(\Gamma):\Gamma\text{ a component of }Q\}
       \leq k-1.
\]
Choose a Kähler class \(\eta\) on \(Y\) (an ample class when \(Y\)
is projective) and a Kähler form \(\omega\) on \(X\).  On
\(H^{1,1}(X,\R)\) define
\[
 \langle \alpha,\beta\rangle_l
       =\int_X \alpha\beta(b^*\eta)^l\omega^{\,n-l-2}.
\]
In this pairing a divisor or rational line denotes its cohomology class.
The exponent is nonnegative since \(l\leq k-1\leq n-2\).
The restriction of a modified-nef class to a resolution of every
prime is pseudo-effective.  Applying this to \(\{P_0+Q\}\)
and pairing on each component of \(Q\) gives
\(\langle P_0+Q,\Gamma\rangle_l\geq0\).
The \(P_0\)-term vanishes by dimension, since it is a pullback
from \(Y\) and \(\dim b(\Gamma)\leq l\).  Summing with the
coefficients of \(Q\) gives
\begin{equation}
                         \langle Q,Q\rangle_l\geq0.
\label{rank:nonnegative-square}
\end{equation}

We now separate images of codimension at least two, where the mixed Hodge
index theorem gives strict negativity, from divisorial images, where only
a multiple of the whole fiber can have square zero.

If \(l<k-1\), the mixed Hodge index theorem contradicts this
inequality.  In this form \(b^*\eta\) has positive square because
\(l+2\leq k\), while
\(\langle Q,b^*\eta\rangle_l=0\) by image dimension.  The form is
a limit of the mixed Hodge index forms obtained by replacing
\(b^*\eta\) with \(b^*\eta+\epsilon[\omega]\); it therefore has
at most one positive direction.  In the orthogonal complement
of a positive-square vector it is negative semidefinite, with
equality only in the radical of the whole form.  But
\(\langle Q,\omega\rangle_l>0\), since at least one component of
\(Q\) has image dimension \(l\).  Thus \(Q\) is not in the radical,
and its square is strictly negative, contrary to
Equation~\eqref{rank:nonnegative-square}.

It remains to treat \(l=k-1\).  Terms from primes whose image has
smaller dimension vanish in this form, as do cross terms from
different divisorial images.  Fix a prime divisor \(D\subset Y\)
and take all the primes \(E_1,\ldots,E_s\subset X\) dominating it.
Write \(a_i=\ord_{E_i}(b^*D)>0\) and
\(C_{ij}=\langle E_i,E_j\rangle_{k-1}\).
The divisor \(D\) is rational Cartier, also when \(Y=W_m\).
Consequently
\[
       C_{ij}\geq0\ (i\ne j),\qquad
       \sum_j C_{ij}a_j=0.
\]
The first assertion is positivity of the proper intersection of
distinct effective primes.  The second pairs \(E_i\) with the
pulled-back divisor \(D\); the additional base factor makes the
intersection vanish by dimension.  Components of \(b^*D\)
mapping into a proper subset of \(D\) make no contribution.
The graph with edges \(C_{ij}>0\) is connected.  Indeed over a
very general point of \(D\) the fiber is connected, all its points
lie in the components \(E_i\), and an intersection contributes a
positive entry precisely when it dominates \(D\).

For any real vector \(x=(x_i)\), the two displayed properties give
the exact identity
\[
 x^{\mathsf T}Cx
   =-\sum_{i<j}C_{ij}a_i a_j
             \left(\frac{x_i}{a_i}-\frac{x_j}{a_j}\right)^2.
\]
Thus the block is negative semidefinite with kernel spanned by
\((a_i)\).  Equation~\eqref{rank:nonnegative-square} forces every
nonzero block of \(Q\) to equal a positive scalar multiple of
\((a_i)\).  In particular \(Q\) then contains every prime over
that divisorial image \(D\).

In the second case this is impossible.  The morphism
\(\nu:W\to W_m\) is an isomorphism over the generic point of
\(D\), and \(E_W\) is exceptional.
Equation~\eqref{rank:zero-minimum} therefore supplies a component over
\(D\) whose \(A\)-coefficient is zero.  Its \(Q\)-coefficient is
zero as well, contradicting the preceding conclusion.

In the first case, the same zero minimum shows that each
divisorial image of a nonzero block of \(Q\) is a component of
\(E_W=N(H)\).  Here \(\{Q\}\) itself is modified nef.  The class
\[
                         b_*(\{Q\}[\omega]^{\,n-k})
\]
is also modified nef on the smooth \(W\).  To see this directly,
choose small-perturbation positive currents for \(\{Q\}\) with
zero generic divisorial Lelong numbers, wedge them with
\(\omega^{n-k}\), and push forward.  Such a current has no trace
mass on a divisor, so its push has no mass on a base divisor:
the inverse image is a union of divisors and subsets of higher
codimension, all of zero trace mass.  The pushed currents
therefore have zero generic divisorial Lelong numbers.  Their
classes tend to the displayed class; a fixed smooth representative
for the vanishing error converts this into the defining
small-Kähler-perturbation condition for modified nefness.
On the other hand, pushing the actual divisor \(Q\) in this
formula kills the components with smaller image and gives a
nonzero positive linear combination of the components of \(N(H)\).
Boucksom's exceptionality of the components of a divisorial negative
part says that no such combination is modified nef
\cite[Definition~3.10 and Theorem~3.12]{Boucksom2004}.
This final contradiction proves \(Q=0\),
and hence Equation~\eqref{rank:upstairs-decomposition}.
\end{proof}

\subsection{Generation from a big nef line}

The remaining nef part is pulled back from a big nef line on a
projective variety.  We use generation along the log canonical boundary
to prove that the adjoint itself is semiample.

\begin{proposition}[Generation from a big nef base line]
\label{rank:big-nef-generation}
Assume \(\Ggood_d\) for \(d<n\).  Let \((Z,\Delta_Z)\) be a
globally \(\Q\)-factorial dlt compact Kähler pair of dimension \(n\)
satisfying the resolution condition of
Definition~\ref{def:lc-strata-resolution}, with effective rational boundary
and analytically nef adjoint
\(L=K_Z+\Delta_Z\).  Suppose there are a smooth compact Kähler
manifold \(U\), a projective resolution \(r:U\to Z\), a proper
surjection \(f:U\to Y\) to a normal projective variety, and a
big nef rational line \(H_Y\) on \(Y\) such that
\begin{equation}
                        r^*L\sim_{\Q}f^*H_Y
\label{rank:big-pullback}
\end{equation}
as actual rational lines.  Then \(L\) is semiample.
\end{proposition}

\begin{proof}
Put \(D_Z=\lfloor\Delta_Z\rfloor\).  We first prove the following assertion
for every pair and diagram satisfying the hypotheses of this proposition:
\begin{equation}
                     \mathbf B(L)\cap D_Z=\varnothing.
\label{rank:floor-avoidance}
\end{equation}
Here \(\mathbf B(L)\) is the intersection of the base loci of all positive
Cartier multiples of \(L\).  Once this is proved, a log canonical threshold
at any remaining stable base locus will create a new floor there and give a
contradiction.

Take a higher log resolution if needed and put
\[
 K_U+F=r^*L=:\ell,\qquad C=F^{=1},\qquad
 E=C-\lfloor F\rfloor,\qquad \Delta_F=F-\lfloor F\rfloor.
\]
The equality defining \(F\) uses the canonical meromorphic
identification.  The divisor \(F\) is an SNC subboundary.
The integral divisor \(E\) is effective, \(r\)-exceptional, and
has no component in common with \(C\): a nonexceptional
coefficient of \(F\) is a coefficient of the effective
\(\Delta_Z\), while every coefficient of \(F\) is at most one.
The image of \(C\) is contained in \(D_Z\).  By
Theorem~\ref{floor:generation}, a divisible multiple of
\(L|_{D_Z}\) is globally generated on the whole reduced floor.

Suppose first that a component \(C_i\) dominates \(Y\).
Pulling the floor sections to \(C_i\) shows that
\((f|_{C_i})^*H_Y\sim_{\Q}\ell|_{C_i}\) is semiample.  Semiampleness
descends under a proper surjection to a normal space in this
situation.  Factor \(C_i\to Y\) through its normal Stein factor
\(Y_i\).  Projection formula descends generation
through the connected-fiber map.  For the finite map \(Y_i\to Y\),
at each \(y\in Y\) choose a section of the generated pullback line
avoiding every point of the finite fiber; finitely many evaluation
kernels cannot cover the section space.  Lemma~\ref{bir:finite-norm}
gives a norm section of a fixed power of the line on \(Y\), nonzero
at \(y\).  These sections generate a fixed multiple of \(H_Y\).
Equation~\eqref{rank:big-pullback}
and normal birational descent then give semiampleness of \(L\).

We may therefore assume that \(C\) is vertical over \(Y\).
Choose a sufficiently divisible integer \(u>1\) for which
\(uL|_{D_Z}\) is generated, and consider the integral line
\begin{equation}
 \mathscr N
   :=\OO_U(u\ell+E-C)
    =\OO_U\bigl(K_U+\Delta_F+(u-1)\ell\bigr).
\label{rank:vanishing-line}
\end{equation}
Although the expression on the right uses rational divisors, its
sum is the integral line on the left.  The restriction sequence is
\[
 0\longrightarrow\mathscr N
 \longrightarrow\OO_U(u\ell+E)
 \longrightarrow\OO_C(u\ell+E)\longrightarrow0.
\]
The direct image of the last sheaf is supported on the proper analytic
subset \(f(C)\), so it is torsion.  Consequently, the two assertions
\begin{equation}
 R^1f_*\mathscr N\ \text{torsion-free},
             \qquad H^1(Y,f_*\mathscr N)=0
\label{rank:vanishing}
\end{equation}
will give surjectivity on global sections: torsion-freeness makes the
connecting map to \(R^1f_*\mathscr N\) zero, and the \(H^1\)-vanishing
then removes the obstruction to lifting a global section.

We prove Equation~\eqref{rank:vanishing} using analytic injectivity.
Equation~\eqref{rank:big-pullback} identifies
\(\mathscr N\otimes K_U^{-1}\) with
\(\Delta_F+(u-1)f^*H_Y\) in \(\Pic(U)\otimes\Q\); this comparison
need not be an isomorphism of the corresponding integral lines.  Fix a
very ample line \(A_Y\) on \(Y\) and its Fubini--Study metric.
A Kodaira decomposition of the big nef \(H_Y\), given arbitrarily
small weight as in the proof of Lemma~\ref{rank:lift}, writes
\((u-1)H_Y\) as a positive rational multiple of \(A_Y\) plus
an effective rational divisor whose fixed singular contribution
is arbitrarily small and whose remaining contribution is a
general divided ample member.  Choose the small weight below
the log canonical threshold of its pullback relative to the
klt SNC boundary \(\Delta_F\), and choose the ample member generally.
Then the resulting boundary on \(U\) is klt.  Choose a positive
integer \(q_0\) clearing all its denominators and realizing both the
chosen Kodaira decomposition and the preceding rational-line comparison
as line isomorphisms, in particular
\[
 (\mathscr N\otimes K_U^{-1})^{\otimes q_0}
 \simeq \OO_U(q_0\Delta_F)\otimes
             f^*\OO_Y\bigl(q_0(u-1)H_Y\bigr).
\]
The divisor metrics and the Fubini--Study metric
define a singular Hermitian metric on
\((\mathscr N\otimes K_U^{-1})^{\otimes q_0}\).  Its \(q_0\)-th
root is a metric on the actual line \(\mathscr N\otimes K_U^{-1}\),
with multiplier ideal \(\OO_U\) and curvature dominating a positive
multiple of the pullback Fubini--Study form.  The same statements hold after every
nonnegative twist by \(f^*A_Y\).

The injectivity theorem \cite[Theorem~A]{FujinoMatsumura} applies
on the compact Kähler \(U\): multiplication by the pullback of
any nonzero section of a positive power of \(A_Y\) is injective
on
\[
             H^1(U,\mathscr N\otimes f^*A_Y^j)
             \quad(j\geq0)
\]
into the correspondingly higher twist.  Its assumptions are
exactly the semipositive smooth metric on the multiplying line,
the preceding positive lower curvature bound, and the trivial
multiplier ideal.
This injectivity implies Equation~\eqref{rank:vanishing}.
For the second assertion, the Leray edge map injects
\(H^1(Y,f_*\mathscr N)\) into \(H^1(U,\mathscr N)\).
Multiplication by a section of a sufficiently high \(A_Y\)-power
sends this subspace to zero, since Serre vanishing kills
\(H^1(Y,f_*\mathscr N\otimes A_Y^j)\) for large \(j\).
Injectivity upstairs forces the original subspace to be zero.
For the first assertion, suppose \(R^1f_*\mathscr N\) has a
nonzero torsion subsheaf.  A section of a sufficiently high
power of \(A_Y\) annihilates a nonzero such subsheaf.
After a further sufficiently large twist it has a nonzero
global section, and Serre vanishing and Leray identify that
section with a nonzero class in
\(H^1(U,\mathscr N\otimes f^*A_Y^j)\).
Multiplication annihilates this class, again contradicting
injectivity.  This proves both assertions.

The resulting surjectivity extends the floor sections as follows.
Pull any section of \(uL|_{D_Z}\)
to \(C\), multiply by the canonical section of \(E|_C\), and
extend by this surjectivity.  The extended section descends to
\(Z\), since \(r_*\OO_U(E)=\OO_Z\) by exceptionality and
normality.  On the strict transforms of the components of
\(D_Z\), division by the canonical section of \(E\) recovers
the prescribed section.  The descended restriction therefore
equals it on the whole reduced \(D_Z\), since equality holds
generically on every component.  Since \(uL|_{D_Z}\) is generated,
this proves Equation~\eqref{rank:floor-avoidance}.

We finish by showing that this stable base locus is empty.  The
base loci of factorial Cartier multiples form a descending
sequence of compact analytic subsets and hence stabilize.
Sections exist because \(H_Y\) is big and
Equation~\eqref{rank:big-pullback} descends their pullbacks.
Choose a divisible degree \(v\) whose base locus is
\(\mathbf B(L)\), and let \(\mathfrak b\) be its base ideal.
If \(\mathbf B(L)\ne\varnothing\), the ideal is a unit near
\(D_Z\), while \((Z,\Delta_Z)\) is klt outside \(D_Z\).
Its log canonical threshold
\[
             t=\operatorname{lct}_{(Z,\Delta_Z)}(\mathfrak b)
\]
is therefore a positive rational number, attained by a divisor
whose center is contained in \(\mathbf B(L)\).
This follows directly on a simultaneous log resolution of the
pair and the ideal: the threshold is the minimum of the positive
rational discrepancies divided by the positive integral ideal
orders.

Choose an integer \(s>t\) and general divisors
\(D_1,\ldots,D_s\in|vL|\).  On the same resolution their fixed
part is the divisor of \(\mathfrak b\); their free transforms
are jointly transverse to the SNC data.  Hence
\[
 \left(Z,\ \Delta_Z+\frac{t}{s}\sum_{i=1}^s D_i\right)
\]
is log canonical: the fixed part is at its log canonical
threshold and the free components have coefficients \(t/s<1\).
There is a new lc place centered in \(\mathbf B(L)\), and the
new adjoint satisfies the actual identity
\[
             L_{\mathrm{new}}\sim_{\Q}(1+tv)L.
\]
Apply \cite[Theorem~3.3]{HLX} to this compact Kähler lc rational
pair.  It gives a projective strongly \(\Q\)-factorial dlt
modification \(d:Z'\to Z\).  Since the pair is lc, all extracted
divisors have log discrepancy zero in our convention, so the
modification is crepant.  The source is again compact Kähler.
Choose its defining dlt log resolution, with exceptional crepant
coefficients strictly below one and an isomorphism at the general
point of every lc stratum; resolving the remaining data away from
those points gives the resolution required by
Definition~\ref{def:lc-strata-resolution}.  The floor on \(Z'\) has a point over the center
of this lc place, hence over \(\mathbf B(L)\).
The new nef adjoint is \(d^*L_{\mathrm{new}}\).  On a common
resolution with \(U\), it satisfies
Equation~\eqref{rank:big-pullback} with the big nef line
\((1+tv)H_Y\).  The previously proved assertion
\eqref{rank:floor-avoidance} therefore makes its stable base locus
disjoint from its floor.  On the
other hand, normality and projection formula identify all
divisible sections under \(d\), and give
\[
       \mathbf B(d^*L_{\mathrm{new}})
                        =d^{-1}\mathbf B(L).
\]
This locus contains the stated floor point, a contradiction.
Thus \(\mathbf B(L)=\varnothing\); factorial stabilization
provides an actual globally generated multiple of \(L\).
\end{proof}

\subsection{Completion of the rank-one case}

\begin{proof}[Proof of Proposition~\ref{rank:rank-one}]
By Proposition~\ref{bir:resolution-transfer} we can begin on the
resolved source in the statement.  We use a secondary induction on
the positive base dimension \(k\).  Apply
Lemma~\ref{rank:comparison} and then Lemma~\ref{rank:base-psef}
to obtain the prepared data and the pseudo-effective line \(H\).

If \(a(W)=0\), Lemma~\ref{rank:zero-base} and the first case of
Lemma~\ref{rank:lift} give
\(J\sim_{\Q}N(J)\), which is already the required decomposition
with trivial positive part.  If \(W\) is projective, apply
Lemma~\ref{rank:base-reduction}.  Its first alternative has a
strictly smaller positive base dimension, so the secondary induction
and birational transfer finish the proof.  In its second alternative,
take a common smooth base resolution in
Equation~\eqref{rank:base-endpoint} and prepare the source over it.
Equation~\eqref{rank:higher-base} adds an effective divisor
exceptional over the old base.  It is also exceptional over \(W_m\),
because the base program extracted no divisors.  Thus the second
case of Lemma~\ref{rank:lift} applies.

We have obtained
\[
                    J\sim_{\Q}P_0+A,\qquad A=N(J)\geq0,
\]
with \(P_0\) nef.  A rationally trivial \(P_0\) finishes the proof.
In the remaining case \(P_0=b^*H_m\), where \(b:X\to W_m\)
is proper and \(H_m\) is big and nef on the projective \(W_m\).
Apply Proposition~\ref{prog:contract-negative} to this
nef-plus-negative presentation.  It gives a globally strongly
\(\Q\)-factorial dlt compact Kähler model \((Z,\Delta_Z)\) with nef
adjoint \(L_Z\).  The ordinary dlt endpoint has the resolution property
of Definition~\ref{def:lc-strata-resolution}.  On a common smooth resolution
\(\mu:U\to X\), \(r:U\to Z\) its exact positive-line comparison is
\[
                  r^*L_Z\sim_{\Q}\mu^*P_0
                        =(b\mu)^*H_m .
\]
Proposition~\ref{rank:big-nef-generation} makes \(L_Z\) semiample.
Hence \(\mu^*P_0\) is semiample.  Lemma~\ref{bir:pullup} identifies
\(\mu^*A=N(\mu^*J)\), so this is precisely \(\Ggood_n\) on \(U\).
Finally Proposition~\ref{bir:resolution-transfer} transfers the
decomposition back through all the prepared source modifications.
\end{proof}
